% 第 3 章 Good Lambda and Weighted Norm Inequalities (原书第 61-74 页)
\chapter{Good \texorpdfstring{$\lambda$}{lambda} 不等式与加权范数不等式}
\label{chap:3}

上一章中, 我们得到了用 $f$ 的大小估计 $Hf$ 大小的整体性结果. 这还不是这个课题的终点; 我们所做的那几类证明提示, 或许存在用 $f(x)$ 或 $f^{\star}(x)$ 估计 $Hf(x)$ 的办法. 我们也许能得到逐点估计, 而不是 $L^p$ 估计.

我们实际上会做一些接近此事、但更弱的事情. 若 $I$ 是一个区间, 我们能指望 $\int_I|Hf|^p\leq C\int_I|f|^p$ 吗? 设想 $I$ 是一个非常小的区间, 这就接近逐点信息了, 但它仍让我们得以使用函数论与 $L^p$ 空间的想法. 为理解这一点, 我们的模型问题是: 若 $\omega>0$, 是否对所有的 $f$ 都有
\[
\int|H f|^{p}\omega\leq C\int|f|^{p}\omega\,?
\]
这称为一个``加权范数不等式''; 我们可以把 $\omega$ 想成把大部分质量集中在一个区间上、而在外面非常小的权重. 我们将用来分析这个问题的技术工具称为 \emph{good-$\lambda$ 不等式} (good-$\lambda$ inequality). 它把 $Hf$ 的大小与 $f^{\star}$ 的大小作比较, 并对``$f^{\star}$ 控制 $Hf$ 的大小''这一直观给出精确的估计.

\section{Good \texorpdfstring{$\lambda$}{lambda} 不等式}
\label{sec:3.1}
\index{good-$\lambda$ 不等式 (good-$\lambda$ inequality)|emph}

我们的直观是: 当 $Hf$ 很大时, $f^{\star}$ 必然也很大. 下面的定理精确地表达了这一点: $Hf$ 很大而 $f^{\star}$ 很小的那个集合, 其测度可以被做得很小.

\begin{theorem}\label{thm:3.1}
存在 $C>0$, 使得对一切 $\gamma<\frac{1}{2}C$\footnote{译者注: 原书此处印作 $\gamma>\frac{1}{2}C$, 应为 $\gamma<\frac{1}{2}C$, 否则与证明及定理~\ref{thm:3.2} 中 $\gamma$ 的取法矛盾.} 与一切 $\lambda$,
\[
|\{x \mid H^{\star}f(x)>2\lambda,\ f^{\star}<\gamma\lambda\}|\leq C\gamma\,|\{x \mid H^{\star}f(x)>\lambda\}|.
\]
\end{theorem}

\begin{proof}
我们先把问题局部化到区间上; 对集合 $B=\{H^{\star}f>\lambda\}$ 应用 Whitney 分解定理, 即定理~\ref{thm:2.30}. 我们得到互不相交的区间 $I$, 满足 $B=\bigcup I$, $|I|\leq d(B^{c},I)\leq 4|I|$. 由于诸区间互不相交, 只需证明
\[
|\{x\in I \mid H^{\star}f(x)>2\lambda,\ f^{\star}<\gamma\lambda\}|\leq C\gamma|I|.
\]
我们用 $f$ 控制 $H^{\star}f$ 的模型是什么? 在定理~\ref{thm:2.40} 的证明中控制 c) 型的项时, $x$ 靠近区间 $J$ 而 $h_j$ 支在 $J$ 中. 这给了我们用 $f^{\star}$ 控制 $H^{\star}f$ 的办法. 这里 $f$ 并不支在 $I$ 中, 对那些不靠近 $I$ 的 $x$ 还会有额外的项. Whitney 分解替我们控制住了这一点: 若 $x$ 不靠近 $I$, 则 $x$ 靠近 $B^{c}$, 而在 $B^{c}$ 上 $H^{\star}f\leq\lambda$.

这一直观提示把 $f$ 分解为靠近 $I$ 的部分与集中于 $\tilde{I}^{c}$ 的部分. 我们选两个参考点来控制 $f$ 的大小: 取 $z_0\in I$ 使 $f^{\star}(z_0)\leq\gamma\lambda$; 如同~\ref*{thm:2.40}c) 中那样, 它将控制 $H^{\star}f$ 的局部部分. 我们再取 $x_0\in B^{c}$, 并免费得到 $H^{\star}f(x_0)\leq\lambda$. 由于 $d(B^{c},I)\leq 4|I|$, 我们可以选取 $x_0$ 使 $d(x_0,I)\leq 5|I|$, 于是这将确定``靠近 $I$''与``远离 $I$''的含义. 但有技术上的困难, 因为 $H^{\star}f(x)$ 与 $H^{\star}f(x_0)$ 不能直接比较; 后者要求我们用在 $x_{0}$ 处居中的区间来度量, 而 $I$ 并不以 $x_0$ 为中心. 我们构造一个比 $I$ 大、但至少以 $x_{0}$ 为中心的区间 $J$. 开始.

取 $z_0\in I$ 使 $f^{\star}(z_0)<\gamma\lambda$; 取 $x_0\in B^{c}$, $d(x_0,I)\leq 5|I|$. 取一个以 $x_0$ 为中心、长度为 $I$ 的 $20$ 倍的区间 $J$. 显然 $I\subset J$, 因为若 $x\in I$, 则 $|x-x_{0}|\leq d(x_{0},I)+|I|\leq 6|I|$.

令 $g=f\chi_J$, $b=f\chi_{J^{c}}$. 我们将证明
\[
\left|\{x\in I \mid H^{\star}g(x)>\frac{\lambda}{2}\}\right|\leq\frac{C}{\lambda}|I|\,f^{\star}(x_{0})\leq C\gamma|I|;
\]
\[
H^{\star}b(x)\leq H^{\star}f(x_{0})+C f^{\star}(x_{0})\leq\lambda+C\gamma\lambda.
\]
有了这些结果,
\begin{align*}
|\{x\in I \mid H^{\star}g(x)>2\lambda,\ f^{\star}<\gamma\lambda\}|
&\leq|\{x\in I \mid H^{\star}g(x)>2\lambda-H^{\star}b,\ f^{\star}<\gamma\lambda\}| \\
&\leq|\{x\in I \mid H^{\star}g(x)>2\lambda-C\gamma\lambda,\ f^{\star}<\gamma\lambda\}| \\
&\leq\left|\{x\in I \mid H^{\star}g(x)>\frac{\lambda}{2}\}\right|\leq C\gamma|I|,
\end{align*}
至少当 $\gamma<\frac{1}{2}C$ 时如此.

对 $H^{\star}g$ 在 $I$ 上的估计并不困难, 尽管这本该是 $H^{\star}$ 最糟糕的部分, 因为我们已经知道 $H^{\star}$ 是弱 $(1,1)$ 的:
\begin{align*}
\left|\left\{x\in I \mid H^{\star}g(x)>\frac{\lambda}{2}\right\}\right|
&\leq\frac{C}{\lambda}\|g\|_{1}=\frac{C}{\lambda}\int_{J}|f| \\
&=\frac{C}{\lambda}|J|\,|J|^{-1}\int_{J}|f|
=20\,\frac{C}{\lambda}|I|\,\frac{1}{|J|}\int_{J}|f|.
\end{align*}
但 $z_0\in I\subset J$, 于是
\[
\frac{1}{|J|}\int_{J}|f|\leq f^{\star}(z_{0})<\gamma\lambda;
\]
整个这一项被 $20C\gamma|I|$ 控制.

为了用 $H^{\star}f(x_{0})$ 控制 $H^{\star}b(x)$, 我们试图把
\[
\int_{|x-y|>\epsilon}\frac{b(y)}{x-y}\,dy
\]
用
\[
\int_{|x_{0}-y|>\epsilon}\frac{b(y)}{x_{0}-y}\,dy
\]
表示出来. 这引入两类误差: 来自 $\frac{1}{x-y}-\frac{1}{x_0-y}$ 的误差, 与来自 $\{|x_0-y|>\epsilon\}\,\Delta\,\{|x-y|>\epsilon\}$ 的误差. 为控制它们, 令
\begin{align*}
S_{1}&=\{y \mid |x_{0}-y|>\epsilon,\ |x-y|<\epsilon\}, \\
S_{2}&=\{y \mid |x_{0}-y|<\epsilon,\ |x-y|<\epsilon\}, \\
E_{j}&=\int_{S_{j}}\frac{|b(y)|}{|x-y|}\,dy.
\end{align*}
则:
\begin{align*}
\left|\int_{|x-y|>\epsilon}\frac{b(y)}{x-y}\,dy\right|
&\leq\left|\int_{|x_{0}-y|>\epsilon}\frac{b(y)}{x-y}\,dy\right|+E_{1}+E_{2} \\
&\leq\left|\int_{|x_{0}-y|>\epsilon}\frac{b(y)}{x_{0}-y}\,dy\right|
+\int_{|x_{0}-y|>\epsilon}\left|\frac{1}{x_{0}-y}-\frac{1}{x-y}\right|dy+E_{1}+E_{2}.
\end{align*}
如所安排, 第一项被 $H^{\star}b(x_0)$ 控制; 不幸的是, 我们在 $x_0$ 处拥有的控制是关于 $H^{\star}f(x_0)$ 而非 $H^{\star}b(x_0)$ 的. 这意味着还要做更多工作, 也显示了我们对 $J$ 的选择何以派上用场. 若取 $y\in J^{c}$, 则 $|y-x_0|>10|I|$. 于是, 若 $y\in\operatorname{supp}b=J^{c}$,
\[
\{y \mid |x_{0}-y|>\epsilon,\ y\in J^{c}\}=\{y \mid |x_{0}-y|>\max\{\epsilon,\ 10|I|\}\},
\]
且
\begin{align*}
\left|\int_{|x_{0}-y|>\epsilon}\frac{b(y)}{x_{0}-y}\,dy\right|
&=\left|\int_{\{y \mid |x_{0}-y|>\max\{\epsilon,\,10|I|\}\}}\frac{f(y)}{x_{0}-y}\,\chi_{J^{c}}\,dy\right| \\
&\leq\left|\sup_{\delta}\int_{|x_{0}-y|>\delta}\frac{f(y)}{x_{0}-y}\,dy\right|=H^{\star}f(x_{0}).
\end{align*}
由卷积核的光滑性控制的那一项, 是我们以前处理过的; 它被
\[
\int_{|x_{0}-y|>\epsilon}|x_{0}-y|^{-1}|x-y|^{-1}|x_{0}-x|\,|f(y)|\,\chi_{J^{c}}\,dy
\]
控制. 我们需要用 $f^{\star}(z_0)$ 控制这一项, 于是需要把 $x_0$ 与 $z_0$ 联系起来, 并用这个关系替换 $|x_0-y|$. 不存在这样的关系, 恰恰意味着我们自己来构造一个. 令
\[
F_{k}=\{y \mid 2^{k+1}\epsilon>|x_{0}-y|>2^{k}\epsilon\}.
\]
我们将证明: 若 $x\in I$, $y\in J^{c}$, $y\in F_k$, 则:
\begin{gather*}
k\geq\log\Bigl(\frac{10|I|}{\epsilon}\Bigr)=C_{\epsilon};\qquad
|x_{0}-y|^{-1}\leq(2^{k}\epsilon)^{-1}; \\
|x-y|^{-1}\leq(2^{k-1}\epsilon)^{-1};\qquad
|z_{0}-x_{0}|<2^{k+1}\epsilon;\qquad
|x-x_{0}|<6|I|.
\end{gather*}
若能证明这些, 则由光滑性控制的那一项被
\[
\sum_{k=C_{\epsilon}}^{\infty}\frac{1}{2^{k}\epsilon}\int_{F_{k}}|f(y)|\,dy\,\frac{|I|}{2^{k-1}\epsilon}
\]
控制, 进而
\begin{align*}
\sum_{k=C_{\epsilon}}^{\infty}&\frac{10|I|}{\epsilon}\,2^{-(k-1)}\,\frac{1}{2^{k}\epsilon}\int_{|y-z_{0}|\leq 2^{k+2}\epsilon}|f(y)|\,dy \\
&\leq\frac{40|I|}{\epsilon}\,f^{\star}(z_{0})\,2^{-C_{\epsilon}}\leq 400\,f^{\star}(z_{0}).
\end{align*}
如今的估计需要的只是耐心; 我们举几个例子. 首先, $y\in F_k$ 蕴含 $2^{k+1}\epsilon>|x_0-y|$, 于是 $y\in F_k\cap J^{c}$ 蕴含 $k\geq C_{\epsilon}$. 由于 $y\in F_k$, 自动地有
\[
|x_{0}-y|^{-1}\leq(2^{k}\epsilon)^{-1};
\]
为估计 $|x-y|$,
\[
|x-y|\geq|x_{0}-y|-|x-x_{0}|>2^{k}\epsilon-6|I|.
\]
但 $2^{k+1}\epsilon>|x_0-y|>10|I|$, 于是 $6|I|<\frac{6}{10}\cdot 2^{k+1}\epsilon$, 从而 $|x-y|\geq 2^{k-1}\epsilon$. 最后,
\[
|z_{0}-x_{0}|\leq d(x_{0},I)+|I|\leq 6|I|<10|I|<2^{k+1}\epsilon,
\]
且
\[
|x_{0}-x|\leq d(x_{0},I)+|I|\leq 6|I|.
\]
为完成证明, 我们必须用 $f^{\star}(z_0)$ 控制误差项 $E_i$; 我们需要把 $|x_0-y|$ 的信息换成 $|z_0-y|$ 的信息. 对 $E_1$ 项:
\[
S_{1}=\{y \mid |x-y|<\epsilon,\ |x_{0}-y|>\epsilon\}.
\]
由于 $z_0\in I$ 且 $x\in I$, $y$ 应当靠近 $x$, 而 $x$ 到 $z_0$ 的距离至多是 $|I|$. 但这必须远小于 $\epsilon$, 因为 $x\in I$ 而 $y\in J^{c}$. 于是:
\[
\epsilon>|x-y|>10|I|,
\]
\[
|z_{0}-y|<|z_{0}-x|+|x-y|<|I|+\epsilon<2\epsilon,
\]
且
\[
S_{1}\subset\{y \mid |y-z_{0}|<2\epsilon\}.
\]
此外,
\[
|x-y|\geq|x_{0}-y|-|x_{0}-x|>\epsilon-6|I|\geq\epsilon-\frac{6}{10}\epsilon=\frac{2}{5}\epsilon,
\]
于是 $|x-y|^{-1}<\frac{5}{2\epsilon}$, 且
\[
E_{1}\leq\frac{5}{2\epsilon}\int_{|y-z_{0}|<2\epsilon}|f(y)|\,dy\leq 5f^{\star}(z_{0}).
\]
对于来自 $S_2=\{y \mid |x_0-y|<\epsilon,\ |x-y|>\epsilon\}$ 的项, 我们注意到 $y$ 必须靠近 $z_0$, 因为 $x_0$ 靠近它; 同样要把 $|I|$ 与 $\epsilon$ 作比较:
\[
\epsilon>|x_{0}-y|>|x-y|-|x-x_{0}|\geq 20|I|-6|I|=14|I|;
\]
\[
|z_{0}-y|\leq|z_{0}-x_{0}|+|x_{0}-y|\leq 6|I|+\epsilon<2\epsilon.
\]
由于 $|x-y|^{-1}<\epsilon^{-1}$, 故 $E_2<2f^{\star}(z_0)$.
\end{proof}

\begin{theorem}\label{thm:3.2}
若 $0<p<\infty$, 则 $\|H^{\star}f\|_p\leq C_p\|f^{\star}\|_p$.
\end{theorem}

\begin{proof}
\begin{align*}
\|H^{\star}f\|_{p}^{p}
&=p\int_{0}^{\infty}|\{H^{\star}f>\lambda\}|\,\lambda^{p-1}\,d\lambda \\
&=2^{p}p\int_{0}^{\infty}|\{H^{\star}f>2\lambda\}|\,\lambda^{p-1}\,d\lambda \\
&\leq 2^{p}p\int_{0}^{\infty}|\{H^{\star}f>2\lambda,\ f^{\star}<\gamma\lambda\}|\,\lambda^{p-1}\,d\lambda
+2^{p}p\int_{0}^{\infty}|\{f^{\star}>\gamma\lambda\}|\,\lambda^{p-1}\,d\lambda \\
&\leq 2^{p}pC\gamma\int_{0}^{\infty}|\{H^{\star}f>\lambda\}|\,\lambda^{p-1}\,d\lambda
+\frac{2^{p}}{\gamma^{p}}\|f^{\star}\|_{p}^{p}.
\end{align*}
选取 $\gamma$ 使 $1-2^{p}C\gamma>0$, 便得到
\[
\left\|H^{\star}f\right\|_{p}\leq 2\gamma\bigl(1-2^{p}C\gamma\bigr)^{-\frac{1}{p}}\left\|f^{\star}\right\|_{p}. \qedhere
\]
\end{proof}

\section{\texorpdfstring{$A_p$}{Ap} 条件}
\label{sec:3.2}

现在我们在区间上控制了 $H^{\star}f$, 并试图用它来回答这个问题: 对哪些 $\omega>0$, 我们能指望
\[
\int|H f|^{p}\omega\leq C\int|f|^{p}\omega\,?
\]
平凡的条件是 $a<\omega<b$; 既然我们能够指望在区间上控制 $Hf$, 我们可以希望有一个条件说: $\omega$ 与它在区间上的平均不能偏离太多; 当然我们希望这种偏离在所有区间上被一致地控制. 表达这一点的方式有很多; 我们可以指望控制 $|\{x\in I \mid \omega(x)>\omega_I\}|$, 也可以指望同时界住 $\int_I\omega$ 与 $\int_I\omega^{-1}$.

\begin{definition}\label{def:3.3}
\index{$A_p$ 权 (condition $A_p$)|emph}
若 $\omega$ 是局部可积的正函数, 我们说 $\omega\in A_p$, 如果:

当 $p=1$ 时,
\[
\omega_{I}\leq C\inf_{x\in I}\omega(x),
\]
其中 $C$ 与 $I$ 无关.

当 $1<p<\infty$ 时,
\[
\left(\frac{1}{|I|}\int_{I}\omega\right)\left(\frac{1}{|I|}\int_{I}\omega^{\frac{1}{p-1}}\right)^{p-1}\leq C,
\]
其中 $C$ 与 $I$ 无关.

当 $p=\infty$ 时,
\[
\frac{\omega(E)}{\omega(I)}\leq\left(\frac{|E|}{|I|}\right)^{\delta}
\]
对某个与 $I$ 无关的 $\delta$ 与一切可测的 $E\subset I$ 成立.
\end{definition}

\begin{remark}\label{rem:3.4}
若 $\omega\in A_\infty$, 令
\[
E=\{x\in I \mid \omega(x)>\beta\omega_{I}\}.
\]
则 $|E|\leq\beta^{\frac{1}{1-\delta}}|I|$, 于是 $A_\infty$ 条件蕴含 $\omega$ 不偏离其平均太多. 反之, 假设
\[
|\{x\in I \mid \omega(x)>\beta\omega_{I}\}|\leq\beta^{\delta}|I|.
\]
则
\begin{align*}
\int_{I}\omega^{-\frac{1}{p-1}}
&=\frac{1}{p-1}\int\left|\left\{\omega^{-1}\geq\lambda\right\}\right|\lambda^{\frac{1}{p-1}-1}\,d\lambda \\
&\leq\frac{1}{p-1}\int_{0}^{\omega_{I}^{-1}}|I|\,\lambda^{\frac{1}{p-1}-1}\,d\lambda
+p\int_{\omega_{I}^{-1}}^{\infty}\left(\lambda^{-1}\omega_{I}^{-1}\right)^{\delta}|I|\,\lambda^{\frac{1}{p-1}-1}\,d\lambda
=C_{\delta}|I|\,\omega_{I}^{-\frac{1}{p-1}}.
\end{align*}
于是
\[
\left(\frac{1}{|I|}\int_{I}\omega\right)\left(\frac{1}{|I|}\int_{I}\omega^{-\frac{1}{p-1}}\right)^{p-1}\leq C_{\delta}^{p-1},
\]
即 $\omega\in A_p$. 对我们而言, 这个计算意味着 $A_p$ 条件与控制函数对其平均的偏离密切相关.
\end{remark}

\begin{lemma}\label{lem:3.5}
若 $\omega\in A_p$, 则
\[
|\{x\in I \mid \omega(x)\leq\beta\omega_{I}\}|\leq C\beta^{\frac{1}{p-1}}|I|.
\]
\end{lemma}

\begin{proof}
把 $A_{p}$ 条件改写,
\begin{align*}
C|I|&\geq(\omega_{I})^{\frac{1}{p-1}}\int_{I}\omega^{-\frac{1}{p-1}}
\geq(\omega_{I})^{\frac{1}{p-1}}\int_{\{\omega<\beta\omega_{I}\}}\omega^{-\frac{1}{p-1}} \\
&\geq(\omega_{I})^{\frac{1}{p-1}}\int_{\{\omega<\beta\omega_{I}\}}(\beta\omega_{I})^{-\frac{1}{p-1}}
=\beta^{-\frac{1}{p-1}}\left|\left\{x\in I \mid \omega<\beta\omega_{I}\right\}\right|. \qedhere
\end{align*}
\end{proof}

\begin{remark}\label{rem:3.6}
对本质上为常数的函数, 不同 $p$ 的 $L^p$ 范数没有差别. ``$A_p$ 函数与其平均相差不多''这一概念使我们能够对 $A_p$ 函数证明基本相同的结果. 我们将证明 $A_p$ 函数既属于 $A_\infty$ 也属于 $A_{p-\epsilon}$. 这两个事实都来自``$L^p$ 范数等价性''的一个精确表述. 记得我们是对 $\omega$ 在区间上作平均, 所以包含关系
\[
L^{p}\subset L^{p+\epsilon}
\]
由 H\"older 不等式得出. 而 $A_p\subset A_{p-\epsilon}$ 这个令人惊讶的结果来自一个``反 H\"older 不等式''.
\end{remark}

\begin{theorem}\label{thm:3.7}
\index{反 H\"older 不等式 (reverse H\"older inequality)}
设 $\omega\in A_p$ 对某个 $p$ 成立, $1<p<\infty$. 则存在 $a$, $q>1$, 使
\[
\left(\frac{1}{|I|}\int_{I}\omega^{q}\right)^{\frac{1}{q}}\leq C\left(\frac{1}{|I|}\int_{I}\omega\right).
\]
\end{theorem}

\begin{proof}
我们先证明: 对某个 $\beta<1$, 有一个 good-$\lambda$ 不等式:
\[
\omega\{x\in I \mid \omega(x)>\lambda\}\leq C\lambda\left|\left\{x\in I \mid \omega(x)>\beta\lambda\right\}\right|.
\]
假设它成立, 定理就容易了:
\[
\int_{I}\omega^{q}=\beta^{-q}C_{q}\int_{0}^{\infty}\lambda^{q-1}\left|\left\{x\in I \mid \omega>\beta\lambda\right\}\right|d\lambda
\]
\begin{align*}
&\geq C'q\int_{\omega_{I}}^{\infty}\lambda^{q-2}\int_{\{x\in I \mid \omega\geq\lambda\}}\omega(x)\,dx\,d\lambda \\
&=C'q\int_{\{x\in I \mid \omega>\omega_{I}\}}\omega(x)\int_{\omega_{I}}^{\omega}\lambda^{q-2}\,d\lambda\,dx \\
&=C'q\int_{\{x\in I \mid \omega>\omega_{I}\}}\omega(x)\left[\frac{\omega^{q-1}(x)}{q-1}-\frac{\omega_{I}^{q-1}}{q-1}\right]dx.
\end{align*}
现在在集合
\[
\{x\in I \mid \omega>\omega_{I}\}
\]
上,
\[
-\omega\,\omega_{I}^{q-1}>-\omega_{I}^{q},
\]
于是
\begin{align*}
C'q\int_{\{x\in I \mid \omega>\omega_{I}\}}\omega(x)\left[\frac{\omega^{q-1}(x)}{q-1}-\frac{\omega_{I}^{q-1}}{q-1}\right]dx
&\geq\frac{C'q}{q-1}\int_{\{x\in I \mid \omega>\omega_{I}\}}\bigl(\omega^{q}-\omega_{I}^{q}\bigr).
\end{align*}
但在 $\{\omega\leq\omega_I\}$ 上, 被积函数是负的, 于是可以把整个这个集合添加进来而不破坏不等式:
\[
\int_{I}\omega^{q}\geq\frac{C'q}{q-1}\int_{I}\bigl(\omega^{q}-\omega_{I}^{q}\bigr),
\]
从而
\[
\left(\frac{C'q}{q-1}-1\right)\frac{1}{|I|}\int_{I}\omega^{q}\leq\left(\frac{1}{|I|}\int_{I}\omega\right)^{q}.
\]
取 $q$ 充分接近 $1$, 证明即告完成.
\end{proof}

还剩下一个 good-$\lambda$ 不等式要证, 它分析 $\omega$ 在区间上的变化. 我们通过一个 Calder\'on--Zygmund 分解来控制它: 构造集合 $B=\{x\in I \mid \omega(x)>\lambda\}$, 并设 $\{J\}$ 是互不相交的区间 $J\subset I$, 满足 $\lambda\leq\omega_J\leq 2\lambda$. 则 $\bigcup J=B$, 且若 $x\in I-B$, 则 $\omega(x)\leq\lambda$. 令 $\alpha=C\beta^{\frac{1}{p-1}}$, 其中 $C$, $\beta$ 如引理~\ref{lem:3.5} 中所述. 若 $\beta<\beta_0$, 则 $\alpha<1$, 且
\[
\omega\{x\in I \mid \omega>\lambda\}=\sum\omega(J)\leq 2\lambda\sum|J|.
\]
但
\[
|J|=|\{x\in J \mid \omega>\beta\omega_{J}\}|+|\{x\in J \mid \omega\leq\beta\omega_{J}\}|.
\]
由于 $\omega_{J}>\lambda$, 它被
\[
|\{x\in J \mid \omega>\beta\omega_{J}\}|+\alpha|J|
\]
控制. 由此得
\[
|J|\leq(1-\alpha)^{-1}\left|\{x\in J \mid \omega>\beta\lambda\}\right|,
\]
最后得
\begin{align*}
\omega\{x\in I \mid \omega>\lambda\}
&\leq C\lambda\sum\left|\left\{x\in J \mid \omega>\beta\lambda\right\}\right| \\
&=C\lambda\left|\left\{x\in B \mid \omega>\beta\lambda\right\}\right|
\leq C\lambda\left|\left\{x\in I \mid \omega>\beta\lambda\right\}\right|.
\end{align*}

\begin{corollary}\label{cor:3.8}
若 $\omega\in A_p$, 则 $\omega\in A_\infty$.
\end{corollary}

\begin{proof}
如定理~\ref{thm:3.7} 那样选取 $q$. 则:
\begin{align*}
\omega(E)&=\int\chi_{E}\,\omega\leq\|\chi_{E}\|_{q'}\,\Bigl(\int\chi_{E}\,\omega^{q}\Bigr)^{\frac{1}{q}}
\leq|E|^{\frac{1}{q'}}\,|I|^{\frac{1}{q}}\left(\frac{1}{|I|}\int_{I}\omega^{q}\right)^{\frac{1}{q}} \\
&\leq C|E|^{\frac{1}{q'}}\,|I|^{\frac{1}{q}}\,\omega_{I}
=C|E|^{\frac{1}{q'}}\,|I|^{\frac{1}{q}-1}\int_{I}\omega
=C\left(\frac{|E|}{|I|}\right)^{\frac{1}{q'}}\omega(I). \qedhere
\end{align*}
\end{proof}

\begin{corollary}\label{cor:3.9}
若 $\omega\in A_p$, 则存在 $\epsilon>0$ 使 $\omega\in A_{p-\epsilon}$.
\end{corollary}

\begin{proof}
改换记号即知 $\omega\in A_p$ 当且仅当
\[
\omega^{-\frac{1}{p-1}}\in A_{p'}.
\]
把定理~\ref{thm:3.7} 应用于 $A_{p'}$, 存在 $q>1$ 使
\[
\left(\frac{1}{|I|}\int_{I}\omega^{-\frac{q}{p-1}}\right)^{\frac{1}{q}}
\leq C\frac{1}{|I|}\int_{I}\omega^{-\frac{1}{p-1}}
\leq C\left(\frac{1}{|I|}\int_{I}\omega\right)^{-\frac{1}{p-1}},
\]
最后一步用了 $\omega\in A_p$. 总之, $\omega\in A_\gamma$, 其中 $\frac{1}{\gamma-1}=\frac{q}{p-1}$; 由于 $q>1$, 故 $\gamma<p$.
\end{proof}

\section{加权范数不等式}
\label{sec:3.3}

\begin{lemma}\label{lem:3.10}
\index{加权范数不等式 (weighted norm inequality)}
若 $\omega\in A_p$, 则
\[
f_{I}\leq C\left(\frac{1}{\omega(I)}\int_{I}|f|^{p}\omega\right)^{\frac{1}{p}}.
\]
若 $\tilde{I}$ 是 $I$ 的倍缩 (即 $\tilde{I}$ 与 $I$ 同心但长度为其三倍), 则 $\omega(\tilde{I})\leq C\omega(I)$; 又若 $J$ 是与 $I$ 同长度、相邻的区间, 则 $\omega(J)\leq c\,\omega(I)$.
\end{lemma}

\begin{proof}
\begin{align*}
f_{I}&=\frac{1}{|I|}\int_{I}f=\frac{1}{|I|}\int_{I}f\,\omega^{\frac{1}{p}}\,\omega^{-\frac{1}{p}} \\
&\leq\frac{1}{|I|}\left(\int_{I}f^{p}\omega\right)^{\frac{1}{p}}\left(\int_{I}\omega^{-\frac{p'}{p}}\right)^{\frac{1}{p'}}.
\end{align*}
但 $\frac{p'}{p}=\frac{1}{p-1}$, 且第二个积分被
\begin{align*}
\left(\int_{I}\omega^{-\frac{p'}{p}}\right)^{\frac{1}{p'}}
&\leq C|I|^{\frac{1}{p'}}\left(\frac{\omega(I)}{|I|}\right)^{-\frac{1}{p(p-1)}} \\
&=C|I|^{\frac{1}{p'}}|I|^{-\frac{1}{p}}\bigl[\omega(I)\bigr]^{-\frac{1}{p}}
\end{align*}
控制.

为得到关于 $\tilde{I}$ 的结果, 令 $f=\chi_{I}$. 则
\[
\frac{1}{2}=f_{\tilde{I}}\leq C\left(\frac{1}{\omega(\tilde{I})}\int_{\tilde{I}}\chi_{I}\,\omega\right)^{\frac{1}{p}},
\]
即
\[
C\leq\frac{\omega(I)}{\omega(\tilde{I})}.
\]
类似的想法可以处理 $J$.
\end{proof}

\begin{lemma}\label{lem:3.11}
\index{倍增条件 (doubling condition)}
设 $f$ 局部可积, 并取 $\omega\in A_p$. 定义
\[
f^{\star\star}(x)=\sup_{x\in I}\frac{1}{\omega(I)}\int_{I}f(y)\,\omega(y)\,dy.
\]
则对 $1<p<\infty$,
\[
\int|f^{\star\star}|^{p}\,\omega\leq C\int|f|^{p}\,\omega.
\]
\end{lemma}

\begin{proof}
回忆定理~\ref{thm:2.24} 的证明; 与那里一样, 我们得到 $E=\{x \mid f^{\star\star}(x)>\lambda\}$ 的一个区间覆盖, 并可以施以选取过程得到互不相交的区间. 我们需要 $E$ 的测度与互不相交区间的测度之和可比, 但 $E\subset\bigcup I^{\star}$, 故 $\omega(E)\leq\sum\omega(I^{\star})$. 而 $I^{\star}$ 的长度是 $I$ 的五倍, 于是引理~\ref{lem:3.10} 告诉我们
\[
\omega(E)\leq\sum\omega(I^{\star})\leq C\sum\omega(I)\leq C\frac{1}{\lambda}\sum\int_{I}f,
\]
等等. 简言之, 只要测度满足``倍增条件''
\[
\omega(I^{\star})\leq C\omega(I),
\]
我们就能对 Lebesgue 测度之外的测度证明极大定理, 即定理~\ref{thm:2.24}.
\end{proof}

\begin{theorem}\label{thm:3.12}
若 $\omega\in A_p$ 对 $1<p<\infty$ 成立, 则
\[
\int|f^{\star}|^{p}\,\omega\leq C\int|f|^{p}\,\omega.
\]
\end{theorem}

\begin{proof}
由于 $\omega\in A_p$, 由推论~\ref{cor:3.9} 存在 $\gamma<p$ 使 $\omega\in A_\gamma$. 则
\[
f_{I}\leq C\left(\frac{1}{\omega(I)}\int_{I}f^{\gamma}\,\omega\right)^{\frac{1}{\gamma}},
\]
从而
\[
f^{\star}(x)\leq C\left[\left(f^{\gamma}\right)^{\star\star}(x)\right]^{\frac{1}{\gamma}}.
\]
于是
\[
\int|f^{\star}|^{p}\,\omega\leq C\int\bigl|\bigl(f^{\gamma}\bigr)^{\star\star}\bigr|^{\frac{p}{\gamma}}\,\omega.
\]
由于 $\frac{p}{\gamma}>1$, 引理~\ref{lem:3.11} 给出最后的积分被
\[
C\int|f^{\gamma}|^{\frac{p}{\gamma}}\,\omega
\]
控制, 这就证明了结果.
\end{proof}

\begin{theorem}\label{thm:3.13}
若 $\omega\in A_p$, 则
\[
\int|H^{\star}f|^{p}\,\omega\leq C\int|f^{\star}|^{p}\,\omega\leq C\int|f|^{p}\,\omega.
\]
\end{theorem}

\begin{proof}
由于 $\omega\in A_p$, 由推论~\ref{cor:3.8} 有 $\omega\in A_\infty$; 按 $A_\infty$ 条件选取 $\delta$. 对 $\{x \mid H^{\star}f(x)>\lambda\}$ 施以 Whitney 分解, 并令
\[
E=\{x\in I \mid H^{\star}f(x)>2\lambda,\ f^{\star}(x)<\gamma\lambda\}.
\]
则由 good-$\lambda$ 不等式, 即定理~\ref{thm:3.1}, 得 $|E|\leq C\gamma|I|$, 因此
\[
\omega(E)\leq\omega(I)\left[\frac{|E|}{|I|}\right]^{\delta}\leq C\gamma^{\delta}\,\omega(I).
\]
于是,
\[
\omega\{x\in I \mid H^{\star}f(x)>2\lambda,\ f^{\star}(x)<\gamma\lambda\}
\leq C\gamma^{\delta}\,\omega\{x\in I \mid H^{\star}f(x)>\lambda\}.
\]
现在 $L^{p}$ 有界性的证明与定理~\ref{thm:3.2} 的证明完全一样.
\end{proof}

\begin{proposition}\label{prop:3.14}
若 $1<p<\infty$, 且
\[
\int|H f|^{p}\,\omega\leq C\int|f|^{p}\,\omega
\]
对所有的 $f\in L^{p}$ 成立, 则 $\omega\in A_{p}$.
\end{proposition}

\begin{proof}
设 $I$, $J$ 如引理~\ref{lem:3.10} 中那样. 则若 $x\in J$, $y\in I$, 有 $|x-y|\leq 2|I|$, 于是若 $f$ 支在 $I$ 中,
\[
|H f(x)|\,\chi_{J}(x)=\left|\int\frac{f(y)}{x-y}\,dy\right|\chi_{J}(x)\geq(2|I|)^{-1}\int_{I}|f|\,\chi_{J}(x).
\]
则
\begin{align*}
\left(\int_{J}\omega\right)|f_{I}|^{p}
&\leq\int|H f|^{p}\,\omega\,\chi_{J}\leq\int|H f|^{p}\,\omega \\
&\leq C\int|f|^{p}\,\omega=C\int_{I}|f|^{p}\,\omega.
\end{align*}
在特殊情形 $f=\chi_{I}$, $f_{I}=1$, 有
\[
\int_{J}\omega\leq C\int_{I}\omega,
\]
于是利用 $I$ 与 $J$ 的对称性,
\[
\omega(I)\leq C\,\omega(J).
\]
于是对支在 $I$ 中的任意 $f$,
\[
\left(\int_{I}\omega\right)|I|^{-p}\left(\int_{I}\omega^{-\frac{1}{p-1}}\right)^{p}
\leq C\int_{I}\omega^{-\frac{p}{p-1}+1},
\]
等等, 等等.
\end{proof}

\begin{lemma}\label{lem:3.15}
若 $\omega\in A_1$, 则 $\omega\in A_p$.
\end{lemma}

\begin{proof}
\begin{align*}
\left(\frac{1}{|I|}\int_{I}\omega\right)\left(\frac{1}{|I|}\int_{I}\omega^{-\frac{1}{p-1}}\right)^{p-1}
&\leq\omega_{I}\left(\frac{1}{|I|}\int_{I}1\right)\left(\inf_{I}\omega\right)^{-1}
=\frac{\omega_{I}}{\inf_{I}\omega}\leq C. \qedhere
\end{align*}
\end{proof}

\begin{corollary}\label{cor:3.16}
\index{$M_s(g)$ 极大函数 (maximal function $M_s(g)$)}
取 $s>1$, 设 $g\geq 0$ 局部可积; 定义 $M_s(g)=[(g^{s})^{\star}]^{\frac{1}{s}}$. 则 $M_s(g)\in A_1$.
\end{corollary}

\begin{proof}
取定 $I$ 与 $\tilde{I}$; 令 $\omega_1=g^{s}\chi_{\tilde{I}}$ 与 $\omega_2=\omega-\omega_1$. 由于 $\frac{1}{s}<1$, 我们可以用推论~\ref{cor:2.27} 控制 $\omega_1^{\star}$:
\begin{align*}
\frac{1}{|I|}\int\bigl(\omega_{1}^{\star}\bigr)^{\frac{1}{s}}
&\leq C\frac{1}{|I|}\int\bigl(\omega_{1}\bigr)^{\frac{1}{s}} \\
&\leq C\frac{1}{|I|}\int_{I}\bigl(\omega_{1}\bigr)^{\frac{1}{s}}
\leq C\bigl(\omega^{\star}(y)\bigr)^{\frac{1}{s}},
\end{align*}
对任意 $y\in I$ 成立. 对所有 $y\in I$ 取下确界, 得
\[
\left[(\omega_{1}^{\star})^{\frac{1}{s}}\right]_{I}\leq C\left[\inf_{y\in I}\omega^{\star}(y)\right]^{\frac{1}{s}}
=C\inf_{y\in I}\bigl[\omega^{\star}\bigr]^{\frac{1}{s}}.
\]
于是 $\omega_{1}$ 被 $A_{1}$ 条件控制.

为控制 $\omega_2$ 这一部分, 取 $x$, $y\in I$ 并计算
\[
\omega_{2}^{\star}(x)=\sup_{J}\frac{1}{|J|}\int_{J}\omega_{2}.
\]
若它非零, 则区间 $J$ 必须与 $\omega_2$ 的支集相交. 另一方面, $J$ 包含 $x$, 所以 $J$ 的长度至少是 $I$ 的长度. 再者, 由于 $x\in J$, $\tilde{J}$ 包含 $y$, 因此
\begin{align*}
\frac{1}{|J|}\int_{J}\omega_{2}
&\leq 2\,\frac{1}{|\tilde{J}|}\int_{\tilde{J}}\omega_{2}
\leq 2\sup_{J}\frac{1}{|J|}\int_{J}\omega=2\omega^{\star}(y).
\end{align*}
现在对所有包含 $x$ 的 $J$ 取上确界, 得 $\omega_2^{\star}(x)\leq 2\omega^{\star}(y)$. 对所有包含 $y$ 的 $I$ 取下确界, 得
\[
\omega_{2}^{\star}(x)\leq 2\inf_{I}\omega^{\star}(y).
\]
于是
\[
\frac{1}{\left|I\right|}\int_{I}\left(\omega_{2}^{\star}\right)^{\frac{1}{s}}
\leq 2\inf_{I}\left(\omega^{\star}\right)^{\frac{1}{s}}(y). \qedhere
\]
\end{proof}

\section{第 3 章评注}
\label{sec:3.4}

\noindent (3.1):\quad 这个结果归功于 R.~Coifman. 这里的证明取自 Coifman 与 C.~Fefferman 的文章~\cite{coifman-fefferman}; 参见该文的致谢. 这类结果同样可以以类似的证明推广到更高维; 如同极大 Hilbert 变换的情形, 我们使用了 $K$ 的逐点估计, 这些证明并不适用于 H\"ormander 型的核.

\noindent (3.2):\quad 利用 good-$\lambda$ 不等式对 $0<p<\infty$ 得到 $L^{p}$ 估计, 是由 D.~Burkholder 从概率论引入的一种技术, 他对与此相关的结果给出了概率论的证明. 见综述文章~\cite{burkholder}.

\noindent (3.3):\quad $A_p$ 条件由 B.~Muckenhoupt~\cite{muckenhoupt-maximal} 在对极大函数与 Hilbert 变换在加权 $L^p$ 空间上行为的一般研究中发现. 诸如 $|x|^\alpha$ 的特殊情形早已为人所知; 见 Hardy--Littlewood--P\'olya~\cite{hardy-littlewood-polya} 第~IX 章的讨论. 我们的方法与 Muckenhoupt 的不同, 同样取自~\cite{coifman-fefferman}. 关于加权范数不等式的更多问题, 见 Muckenhoupt 的综述文章~\cite{muckenhoupt-classical}, 尽管那里提到的问题已经取得了大量进展. $A_2$ 权的理论在调和函数边值的分析中扮演重要角色; 讨论见 P.~Jones~\cite{jones} 的文章.

\noindent (3.16):\quad 这个结果最初由 A.~C\'ordoba 与 C.~Fefferman~\cite{cordoba-feffermanC-weighted} 证明; 这里的证明取自 R.~Coifman 与 R.~Rochberg~\cite{coifman-rochberg}.
